\title{DeepSeekMath: Pushing the Limits of Mathematical Reasoning in Open Language Models}

\begin{abstract}
Mathematical reasoning represents a formidable obstacle for language models because of its highly structured and intricate demands. In this work, we present DeepSeekMath~7B, an extension of the DeepSeek-Coder-Base-v1.5 7B model, further trained on a corpus comprising 120B math-specific tokens gathered from Common Crawl, in addition to data from natural language and code sources. Remarkably, DeepSeekMath~7B achieves a score of 51.7\% on the challenging, competition-caliber MATH benchmark—without resorting to auxiliary toolkits or ensembling approaches—narrowing the gap with leading proprietary models such as Gemini-Ultra and GPT-4. Moreover, employing self-consistency sampling (64 samples) pushes its MATH benchmark accuracy to 60.9\%. The strong mathematical reasoning exhibited by DeepSeekMath~is driven by two principal components: (1) the effective exploitation of large-scale, publicly available web data through a carefully crafted data filtering and selection workflow, and (2) the application of Group Relative Policy Optimization (GRPO)—a refined form of Proximal Policy Optimization (PPO)—which boosts reasoning abilities and simultaneously streamlines PPO's memory requirements.
\end{abstract}

\section{Introduction}

The advent of large language models (LLMs) has transformed AI methodologies for mathematical reasoning, leading to marked progress on tasks such as the quantitative reasoning benchmark \citep{MATH} and geometry reasoning benchmark \citep{trinh2024solving}. These advances have also established LLMs as valuable collaborators in tackling advanced mathematical problems alongside human experts \citep{tao}. Nevertheless, leading models like GPT-4 \citep{gpt4} and Gemini-Ultra \citep{gemini} remain inaccessible to the public, and currently available open-source alternatives lag well behind in terms of capability.

This paper presents DeepSeekMath, a domain-adapted language model that demonstrates superior mathematical reasoning compared to existing open-source systems and rivals GPT-4’s performance on established academic evaluations. Central to this advancement is the development of the DeepSeekMath Corpus—a comprehensive, high-quality pre-training resource comprising 120B mathematics-focused tokens. The dataset is curated from the Common Crawl (CC) by leveraging a fastText-based classifier \citep{joulin2016fasttext}. Initially, this classifier is trained with examples from OpenWebMath \citep{openwebmath} as positive samples, alongside a diverse array of unrelated web pages serving as negatives. After this, the classifier is used to extract more math-relevant samples from the CC, and these selections undergo further refinement through manual annotation. The classifier is then retrained on this enriched data to achieve higher accuracy. Empirical analysis demonstrates that the corpus maintains a high degree of quality: our foundational model, DeepSeekMath-Base 7B, attains 64.2\% on GSM8K \citep{gsm8k} and 36.2\% on the MATH competition dataset \citep{MATH}, surpassing the results of Minerva 540B \citep{minerva}. Furthermore, as the DeepSeekMath Corpus supports multiple languages, notable progress is observed on Chinese mathematics benchmarks \citep{wei2023cmath,agieval}. We regard our methodologies for constructing mathematical datasets as an initial step for the wider research community, with ample scope for continued enhancement.

DeepSeekMath-Base is built upon DeepSeek-Coder-Base-v1.5 7B \citep{deepseek-coder}, given our findings that models pretrained on code exhibit superior performance for subsequent math-specific training compared to those initialized from general-purpose language models. Additionally, mathematical pretraining not only bolsters the model’s proficiency in mathematical tasks, but also leads to enhancements on general reasoning benchmarks such as MMLU \citep{mmlu} and BBH \citep{bbh}, demonstrating an overall strengthening of broad cognitive abilities.

Following pretraining, DeepSeekMath-Base undergoes specialized mathematical instruction tuning, utilizing datasets featuring chain-of-thought \citep{cot}, program-of-thought \citep{pot,pal}, and tool-augmented reasoning approaches \citep{tora}. The fine-tuned DeepSeekMath-Instruct 7B model subsequently outperforms all existing 7B-scale alternatives and exhibits performance on par with leading instruction-tuned models at the 70B parameter level.

In addition, we propose Group Relative Policy Optimization (GRPO), a modified reinforcement learning (RL) method inspired by Proximal Policy Optimization (PPO) \citep{schulman2017proximal}. GRPO removes the necessity for a critic model and, instead, derives baseline estimates based on group score statistics, resulting in notable reductions in training overhead. Leveraging only a portion of English-language instruction tuning data, GRPO yields significant improvements over the already competitive DeepSeekMath-Instruct model—raising performance both on target domains (e.g., GSM8K: 82.9\% $\rightarrow$ 88.2\%, MATH: 46.8\% $\rightarrow$ 51.7\%) and transfer domains (e.g., CMATH: 84.6\% $\rightarrow$ 88.8\%) during RL optimization. We further present a comprehensive framework to interpret a range of methods including Rejection Sampling Fine-Tuning (RFT) \citep{yuan2023scaling}, Direct Preference Optimization (DPO) \citep{dpo}, PPO, and GRPO. Through this unified framework, we show these approaches can all be regarded as forms of direct or simplified RL strategies. We conduct a series of thorough experimental analyses—spanning comparisons such as online versus offline learning, outcome versus process supervision, and single-step versus iterative RL—to disentangle key contributors within this unified structure. Finally, we elucidate the mechanisms underlying the observed gains from RL on instruction-tuned models and propose future research avenues for advancing RL efficacy in this context.

\subsection{Contributions}
This work advances the field through large-scale math pre-training and a systematic investigation into reinforcement learning.

\textbf{Math Pre-Training at Scale}

\begin{itemize}[topsep=0pt]
    \item Our study demonstrates that Common Crawl, a freely available web resource, offers substantial and relevant material for mathematical applications. Leveraging a carefully crafted data selection process, we create the DeepSeekMath Corpus, which comprises 120B tokens extracted from math-relevant web pages. This dataset surpasses Minerva's math web dataset \citep{minerva} by nearly sevenfold and is nine times larger than OpenWebMath \citep{openwebmath}.

    \item With our pre-trained model, DeepSeekMath-Base 7B, we achieve results on par with Minerva 540B \citep{minerva}, suggesting that scale in parameter count is not the sole driver of mathematical reasoning performance. These outcomes reveal that smaller models, when trained on refined, high-quality data, are also capable of attaining strong mathematical competence.

    \item We present insights from our experiments with math pre-training. We observe that performing code training before math-specific training significantly enhances a model’s proficiency in mathematical problem-solving, regardless of tool use. This contributes evidence to the ongoing debate around the value of code pre-training for reasoning skills, supporting the view that code training positively impacts mathematical reasoning.

    \item In contrast to prevailing practices, our findings indicate that incorporating arXiv papers into training, a strategy often seen in math-related research, yields no meaningful benefit across any of the mathematical benchmarks evaluated in this work.
\end{itemize}

\textbf{Exploration and Analysis of Reinforcement Learning}

\begin{itemize}[topsep=0pt]
    \item We present Group Relative Policy Optimization (GRPO), a novel reinforcement learning algorithm designed for efficiency and effectiveness. Unlike traditional approaches, GRPO eliminates the need for a critic model and instead leverages group scores to estimate the baseline, leading to considerable reductions in computational requirements when compared to Proximal Policy Optimization (PPO).
    \item Our empirical results indicate that GRPO brings notable improvements to our instruction-tuned model, DeepSeekMath-Instruct, relying solely on instruction-tuning datasets. Additionally, we report gains in the model's ability to generalize to out-of-domain tasks throughout reinforcement learning.
    \item We propose a cohesive framework to interpret a variety of methodologies—including RFT, DPO, PPO, and GRPO—and undertake thorough experimentation. Our investigations span topics such as online versus offline training, outcome-based versus process-based supervision, as well as single-turn versus iterative reinforcement learning, all aimed at disentangling key aspects of this overarching framework.
    \item Leveraging our unified framework, we analyze the underlying mechanisms contributing to the success of reinforcement learning, and outline several promising research avenues for further enhancing reinforcement learning techniques in large language models (LLMs).
\end{itemize}

\subsection{Summary of Evaluations and Metrics}

\begin{itemize}[topsep=0pt]
    \item \textbf{English and Chinese Mathematical Reasoning}: We perform in-depth evaluations of our models on a range of English and Chinese mathematical reasoning tasks, spanning difficulty levels from primary school up to collegiate mathematics. The English assessment suite includes GSM8K \citep{gsm8k}, MATH \citep{MATH}, SAT \citep{llemma}, OCW Courses \citep{minerva}, and MMLU-STEM \citep{mmlu}. For Chinese, our benchmarks encompass MGSM-zh \citep{mgsm}, CMATH \citep{wei2023cmath}, Gaokao-MathCloze \citep{agieval}, and Gaokao-MathQA \citep{agieval}. The evaluations measure models' abilities to deliver complete, self-sufficient textual solutions as well as their proficiency in computational problem solving using Python.

    In terms of English mathematical tasks, DeepSeekMath-Base demonstrates performance on par with the proprietary Minerva 540B \citep{minerva}, and achieves substantially better results than all other open-source foundational models—such as Mistral 7B \citep{mistral} and Llemma-34B \citep{llemma}—regardless of whether these models have undergone mathematical pre-training. Noteworthy is DeepSeekMath-Base's leading performance on Chinese datasets, a result attributable to the inclusion of diverse, high-quality non-English mathematical data, rather than adhering to the English-centric pre-training methodology seen in earlier work \citep{minerva,llemma}. When enhanced through instruction tuning and reinforcement learning, DeepSeekMath-Instruct and DeepSeekMath-RL reach new benchmarks, with DeepSeekMath-RL surpassing 50\% accuracy on the competitive MATH dataset—a first for open-source models.
\end{itemize}

\item \textbf{Formal Mathematics}: To assess DeepSeekMath-Base, we employ the informal-to-formal theorem proving benchmark introduced in \citep{dsp_proof}, utilizing miniF2F \citep{minif2f} and the Isabelle proof assistant \citep{isabelle}. Results show that DeepSeekMath-Base achieves robust few-shot performance in automatic formalization tasks.
\item \textbf{Natural Language Understanding, Reasoning, and Code}: To thoroughly evaluate the model's competence in general comprehension, logical reasoning, and programming, DeepSeekMath-Base is tested on several comprehensive benchmarks. These include the Massive Multitask Language Understanding (MMLU) dataset \citep{mmlu} covering 57 varied multiple-choice domains, BIG-Bench Hard (BBH) \citep{bbh} which features 23 particularly demanding multi-step reasoning tasks, and the HumanEval \citep{codex} and MBPP \citep{mbpp} benchmarks widely adopted for code generation assessment. Pre-training with mathematical content has been found to enhance the model's proficiency in both understanding language and executing complex reasoning.

\section{Math Pre-Training}

\subsection{Data Collection and Decontamination} This section details the methodology for building the DeepSeekMath Corpus from Common Crawl data. As illustrated in Figure \ref{fig:data_collect}, we describe an iterative framework that facilitates the large-scale extraction of mathematical content from Common Crawl, beginning from an initial seed dataset—such as a compact, carefully curated set of math-specific texts. Notably, this pipeline is adaptable for other specialized domains, like programming.

We commence by designating OpenWebMath \citep{openwebmath}, a vetted repository of mathematical webpages, as our seed corpus. This seed enables us to train a fastText model \citep{joulin2016fasttext} designed to identify web content analogous to OpenWebMath. For this purpose, we randomly sample 500,000 items from the seed as positive examples, and select an equal number of unrelated pages from Common Crawl to serve as negative samples. Utilizing an open-source fastText implementation\footnote{\url{https://fasttext.cc}}, we set parameters as follows: vector dimensionality of 256, learning rate of 0.1, a word n-gram length up to 3, minimum word occurrence threshold of 3, and train for 3 epochs. Prior to mining, Common Crawl undergoes URL-level deduplication and near-duplicate filtering, reducing the dataset to 40 billion distinct HTML documents. Next, we deploy the trained fastText model to extract mathematics-focused pages from this deduplicated web archive. Pages are ranked by their fastText model scores to discard low-quality content, preserving only the highest scoring documents. The impact of dataset size is subsequently assessed by pre-training models on data comprising the top 40B, 80B, 120B, and 160B tokens; for the initial cycle, the 40B token subset is selected for retention.

Following the initial round of data gathering, a considerable number of mathematical web pages are still missing from the collection. This shortfall is mainly attributable to the fastText model’s training on a set of positive samples that does not capture sufficient variety. To enhance the seed dataset, we expand our pool of mathematical web sources, thus aiming to better the performance of the fastText classifier. In detail, we begin by segmenting the entire Common Crawl dataset into distinct domains, where a domain is defined by all pages sharing the same base URL. For every domain, we assess what proportion of web pages were captured in the first iteration. If more than 10\% of a domain’s pages were collected, we categorize that domain as math-oriented (e.g., \url{mathoverflow.net}). 

Next, within these domains, we manually label the URLs that correspond to mathematical content (such as \url{mathoverflow.net/questions}). Web pages connected to these URLs but missed in previous rounds are subsequently incorporated into the seed data. Through this iterative process, we accumulate a richer set of positive examples and, in turn, train an improved fastText model that demonstrates higher recall for mathematical material during the following round. Ultimately, after four such cycles, we obtain 35.5M mathematical pages, comprising a total of 120B tokens. By the fourth pass, nearly 98\% of data had already been collected in the third, prompting us to discontinue further collection.

To mitigate the risk of benchmark leakage, we adopt the methodology from \cite{deepseek-coder}, filtering out web pages that include any questions or solutions from established English mathematical benchmarks like GSM8K \citep{gsm8k}, MATH \citep{MATH}, as well as Chinese datasets such as CMATH \citep{wei2023cmath} and AGIEval \citep{agieval}. The exclusion rule is as follows: any snippet containing a 10-gram that exactly matches a substring in the evaluation sets is omitted from the mathematical training corpus. For reference materials shorter than 10-grams but at least 3-grams, we similarly perform exact matching to exclude compromised data.

\subsection{Validating the Quality of the DeepSeekMath~Corpus}

To assess the DeepSeekMath~Corpus’s relative quality, we conduct pre-training experiments, benchmarking against other contemporary mathematical corpora:

\begin{itemize}[topsep=0pt]
    \item \textbf{MathPile} \citep{mathpile}: This composite corpus (8.9B tokens) includes materials from textbooks, Wikipedia, ProofWiki, CommonCrawl, StackExchange, and arXiv, with over 85\% derived from arXiv;
    \item \textbf{OpenWebMath} \citep{openwebmath}: A selection of CommonCrawl filtered for mathematical relevance, consisting of 13.6B tokens;
    \item \textbf{Proof-Pile-2} \citep{llemma}: A collection merging OpenWebMath, AlgebraicStack (10.3B tokens of mathematical code), and arXiv (28.0B tokens); in line with \cite{llemma}, experiments employ an arXiv:Web:Code token ratio of 2:4:1.
\end{itemize}

\subsubsection{Training Setting}
\label{sec:quality-policy}
We fine-tune a general-purpose pre-trained language model comprising 1.3B parameters, utilizing the same architecture as the DeepSeek LLMs \citep{deepseek-llm}, which we refer to as DeepSeek-LLM 1.3B. Separate models are trained on each mathematical dataset, each for a total of 150B tokens. For all experiments, we employ the lightweight and efficient HAI-LLM training infrastructure \citep{haillm}. In line with the methodologies established for DeepSeek LLMs, the AdamW optimizer \citep{adamW} is utilized with hyperparameters $\beta_1=0.9$, $\beta_2=0.95$, and $\mathrm{weight\_decay}=0.1$. The learning rate follows a multi-phase schedule: it linearly increases and peaks after 2,000 warmup iterations, is reduced to 31.6\% of the maximum after 80\% of training, and further reduced to 10.0\% of the maximum at 90\% completion of training. We set the peak learning rate to 5.3e-4, use a batch size of 4M tokens, and employ a context length of 4K.

\subsubsection{Evaluation Results}

\textbf{The DeepSeekMath~Corpus stands out for its exceptional quality, multilingual composition, and unparalleled scale among available mathematical corpora.}

\begin{itemize}[topsep=0pt]
    \item \textbf{High-quality}:
    Model performance is assessed on eight downstream mathematical tasks, leveraging few-shot chain-of-thought prompting \cite{cot}. Table \ref{tab:corpora_comparison} reveals a distinct advantage for models trained on the DeepSeekMath~Corpus. Furthermore, as illustrated in Figure \ref{fig:corpora_comparisons}, the model trained with DeepSeekMath outperforms that trained on Proof-Pile-2 after processing 50B tokens (representing one full epoch over Proof-Pile-2), suggesting that the DeepSeekMath Corpus possesses a higher average data quality.
    \item \textbf{Multilingual}:
    The DeepSeekMath~Corpus features content in multiple languages, with English and Chinese as the most prominent. Results in Table \ref{tab:corpora_comparison} show that models trained on DeepSeekMath exhibit improved mathematical reasoning abilities in both English and Chinese. By contrast, most existing mathematical datasets are heavily focused on English, offering only marginal benefits—or even impairing—performance on Chinese mathematical problems.
    \item \textbf{Large-scale}:
    In terms of scale, DeepSeekMath~Corpus vastly exceeds other available mathematical corpora. Figure \ref{fig:corpora_comparisons} demonstrates that training DeepSeek-LLM 1.3B on DeepSeekMath results in a steeper and more sustained learning trajectory. Conversely, due to their relatively limited size, existing baseline corpora must be cycled through repeatedly during training, causing their corresponding models’ performance to plateau rapidly.
\end{itemize}

\subsection{Training and Evaluating DeepSeekMath-Base 7B}

Here, we detail the development and evaluation of DeepSeekMath-Base 7B, a foundational model engineered to excel at mathematical reasoning tasks. The model initialization employs DeepSeek-Coder-Base-v1.5 7B \citep{deepseek-coder} as its base, followed by a training regimen spanning 500B tokens. The composition of the training dataset is as follows: 56\% from the DeepSeekMath Corpus, 4\% from AlgebraicStack, 10\% from arXiv, 20\% consists of code extracted from Github, and the remaining 10\% comes from Common Crawl natural language texts in both English and Chinese. Training largely follows the methodology outlined in Section \ref{sec:quality-policy}, with two exceptions: the learning rate peak is fixed at 4.2e-4, and training occurs with a batch size of 10M tokens.

To systematically gauge DeepSeekMath-Base 7B’s mathematical proficiency, we perform a thorough evaluation addressing three principal abilities: formulating complete mathematical solutions independently, tackling problems that require tool use, and conducting formalized theorem proving. Additionally, we characterize the model’s broader competencies, which include understanding and reasoning in natural language as well as programming.

\paragraph{Mathematical Problem Solving with Step-by-Step Reasoning}

DeepSeekMath-Base is assessed on its ability to solve mathematics problems via few-shot chain-of-thought prompting \citep{cot}, utilizing eight diverse benchmarks in both English and Chinese. These include datasets focusing on quantitative reasoning (such as GSM8K \citep{gsm8k}, MATH \citep{MATH}, and CMATH \citep{wei2023cmath}) as well as multiple-choice problem sets (like MMLU-STEM \citep{mmlu} and Gaokao-MathQA \citep{agieval}), spanning a broad array of mathematical topics ranging from basic to university-level challenges.

Table \ref{tab:base_math_cot} illustrates that among open-source base models, DeepSeekMath-Base 7B consistently achieves the highest scores across all eight evaluation benchmarks. This comparison encompasses established models such as Mistral 7B \citep{mistral} as well as newer entries like Llemma 34B \citep{llemma}, which has undergone specialized math training with Proof-Pile-2 \citep{llemma}. Particularly impressive is DeepSeekMath-Base 7B's result on the competition-oriented MATH dataset, where it exceeds other open-source base models by an absolute margin greater than 10\%. Furthermore, it surpasses Minerva 540B \citep{minerva}, a closed-source base model with a parameter count 77 times larger—built upon PaLM \citep{palm} and further trained on mathematical corpora.

\paragraph{Mathematical Problem Solving with Tool Use} For evaluating program-aided mathematical reasoning, we employ few-shot program-of-thought prompting \citep{pot,pal} on both GSM8K and MATH datasets. In this approach, models are tasked with crafting Python programs for each problem, making use of libraries such as \textit{math} and \textit{sympy} to execute complex computations. The correctness of the answer is determined by running the generated code and checking its output. Table \ref{tab:base_math_pot_proof} shows that DeepSeekMath-Base 7B achieves superior performance compared to the prior leading model, Llemma 34B.

\paragraph{Formal Mathematics}

Formal proof automation has garnered considerable interest for its contributions to proof reliability and process efficiency. We assess DeepSeekMath-Base 7B using the informal-to-formal proving task described in \citep{dsp_proof}: generating a formal proof given an informal problem statement, its formal version, and an informal proof. The evaluation is conducted on the miniF2F \citep{minif2f} benchmark, which targets Olympiad-style formal mathematics, requiring the generation of Isabelle formal proofs through few-shot prompting. Following the methodology of \cite{dsp_proof}, the model produces a proof sketch, with Sledgehammer \citep{sledgehammer} providing automated assistance to complete the details. As shown in Table \ref{tab:base_math_pot_proof}, DeepSeekMath-Base 7B exhibits robust autoformalization capabilities.

\paragraph{Natural Language Understanding, Reasoning, and Code}

We further assess DeepSeekMath-Base 7B on natural language understanding with MMLU \citep{mmlu}, reasoning with BBH \citep{bbh}, and code generation on HumanEval \citep{codex} as well as MBPP \citep{mbpp}. Table \ref{tab:understanding_reasoning_code} demonstrates marked improvements in MMLU and BBH performance relative to its predecessor, DeepSeek-Coder-Base-v1.5 \citep{deepseek-coder}, evidencing the benefits of mathematical training on tasks involving language comprehension and reasoning. Additionally, the integration of code tokens during further training allows DeepSeekMath-Base 7B to preserve the code generation performance of DeepSeek-Coder-Base-v1.5 on HumanEval and MBPP. Across the three benchmarks covering reasoning and coding, DeepSeekMath-Base 7B also consistently outperforms the generalist Mistral 7B model \citep{mistral}.

\section{Supervised Fine-Tuning}

\subsection{SFT Data Curation}

We assemble an instruction-tuning dataset for mathematics that encompasses both English and Chinese questions, with coverage across a spectrum of mathematical domains and difficulty. Each problem is paired with a solution expressed in one of several reasoning styles, including chain-of-thought (CoT) \citep{cot}, program-of-thought (PoT) \citep{pot,pal}, or tool-integrated formats \citep{tora}. The dataset includes a total of 776K training samples.

\begin{itemize}[topsep=0pt]
    \item \textbf{English mathematical datasets}: GSM8K and MATH datasets are annotated with solutions integrating external tools. Additionally, we incorporate selected problems from MathInstruct \citep{MathInstruct} and use the training split of Lila-OOD \citep{lila}, which offer CoT and PoT style solutions. Our English-language compilation addresses a broad array of topics such as algebra, probability, number theory, calculus, and geometry.
    \item \textbf{Chinese mathematical datasets}: For Chinese, we aggregate K-12 problems distributed over 76 sub-disciplines (e.g., linear equations), each annotated with solutions featuring both CoT reasoning and tool-integrated approaches.
\end{itemize}

\subsection{Training and Evaluating DeepSeekMath-Instruct 7B}

We detail here the development of DeepSeekMath-Instruct 7B, which has been adapted for mathematical instruction following pre-training on DeepSeekMath-Base. To prepare data for training, problems are concatenated at random up to a context limit of 4K tokens. The model is trained across 500 steps using a batch size of 256, and employs a fixed learning rate of 5e-5.

Mathematical proficiency of our models is assessed both in the absence and presence of external tool usage, utilizing four quantitative reasoning benchmarks administered in English and Chinese. Our model’s results are compared with several state-of-the-art contemporaries:

\begin{itemize}[topsep=0pt]
    \item \textbf{Closed-source models}: This group features (1) the GPT series, notably GPT-4 \citep{gpt4} and GPT-4 Code Interpreter~\footnote{\url{https://openai.com/blog/chatgpt-plugins\#\#code-interpreter}}, (2) Gemini Ultra and Pro \citep{gemini}, (3) Inflection-2 \citep{inflection-2}, (4) Grok-1~\footnote{\url{https://x.ai/model-card}}, along with recent Chinese models such as (5) Baichuan-3~\footnote{\url{https://www.baichuan-ai.com}}, and (6) the most recent GLM-4~\footnote{\url{https://open.bigmodel.cn/dev/api\#glm-4}} from the GLM series \citep{glm}.
\end{itemize}

These models serve general purposes and most have been subjected to various alignment procedures.
    \item \textbf{Open-source models} comprise a range of options: general-purpose systems such as (1) DeepSeek-LLM-Chat 67B \citep{deepseek-llm}, (2) Qwen 72B \citep{qwen}, (3) SeaLLM-v2 7B \citep{seallm}, and (4) ChatGLM3 6B \citep{chatglm3}, alongside specialized models with mathematical enhancements. These include (5) InternLM2-Math 20B~\footnote{\url{https://github.com/InternLM/InternLM-Math}}, an extension of InternLM2 further trained on mathematical data and then subject to instruction tuning, (6) Math-Shepherd-Mistral 7B, which applies PPO training \citep{schulman2017proximal} to Mistral 7B \citep{mistral} utilizing a process-supervised reward model, (7) the WizardMath family \citep{wizardmath}, which boosts mathematical reasoning skills in Mistral 7B and Llama-2 70B \citep{llama2} through evolve-instruct (a variant of instruction tuning relying on AI-generated instructions) and PPO optimization using predominantly GSM8K and MATH for training, (8) MetaMath 70B \citep{metamath}, where Llama-2 70B is fine-tuned on an enriched set of GSM8K and MATH data, (9) ToRA 34B \cite{tora}, based on CodeLlama 34B and adapted for tool-integrated mathematical reasoning, and (10) MAmmoTH 70B \citep{MathInstruct}, where Llama-2 70B undergoes instruction tuning on the MathInstruct dataset.
\end{itemize}

Referring to Table \ref{tab:sft_rl_math}, when assessed under the constraint that prohibits tool usage, DeepSeekMath-Instruct 7B displays exceptional capabilities in stepwise reasoning. Specifically, for the challenging MATH dataset, this model outperforms every open-source alternative and the vast majority of closed-source systems (including Inflection-2 and Gemini Pro), achieving an advantage of at least 9\% in absolute terms. This superiority holds even compared to models with much larger parameter counts (such as Qwen 72B) or those trained with specialized math-focused reinforcement learning protocols (e.g., WizardMath-v1.1 7B). Although DeepSeekMath-Instruct is competitive with Chinese proprietary models like GLM-4 and Baichuan-3 on MATH, it still falls short compared to GPT-4 and Gemini Ultra.

When evaluation permits both natural language reasoning and the application of programmatic tools, DeepSeekMath-Instruct 7B reaches an accuracy of nearly 60\% on MATH, outperforming all prior open-source models. Across other benchmarks, its performance is on par with DeepSeek-LLM-Chat 67B, which previously represented the leading open-source result despite being ten times larger in scale.
\section{Reinforcement Learning}

\subsection{Group Relative Policy Optimization} Reinforcement learning (RL) methods have been shown to effectively bolster the mathematical reasoning performance of LLMs beyond what is achieved with Supervised Fine-Tuning (SFT) \citep{wang2023math,wizardmath}. Here, we present Group Relative Policy Optimization (GRPO), an efficient and robust RL approach.

\subsubsection{From PPO to GRPO} Proximal Policy Optimization (PPO) \citep{schulman2017proximal} remains the prevalent actor-critic algorithm for RL-based fine-tuning of LLMs \citep{ouyang2022training}. This method updates LLM parameters by maximizing the surrogate objective given by:

\begin{equation}
\footnotesize
    \mathcal{J}_{PPO}(\theta) = \mathbb{E}{[q \sim P(Q), o \sim \pi_{\theta_{old}}(O|q)]} \frac{1}{|o|} \sum_{t=1}^{|o|} \min \left[ \frac{\pi_\theta(o_{t} | q, o_{<t})}{\pi_{\theta_{old}}(o_{t} | q, o_{<t})} A_{t}, \text{clip} \left( \frac{\pi_\theta(o_{t} | q, o_{<t})}{\pi_{\theta_{old}}(o_{t} | q, o_{<t})}, 1 - \varepsilon,

Here, $\pi_{\theta}$ and $\pi_{\theta_{old}}$ refer to the current and previous policy networks, respectively, while $q$ and $o$ denote questions and their corresponding generated responses, with $o$ sampled both from the dataset of questions and via the old policy $\pi_{\theta_{old}}$. The hyper-parameter $\varepsilon$ is associated with clipping in PPO, serving to enhance training stability. $A_t$ represents the advantage, determined using Generalized Advantage Estimation (GAE) \citep{gae} and based on the set of rewards $\{r_{\ge t}\}$ along with an estimated value function $V_{\psi}$. Consequently, PPO necessitates joint training of a value function alongside the policy model. In order to prevent excessive reward model optimization, it is standard practice to incorporate a per-token Kullback-Leibler (KL) penalty from a reference policy into the reward function for each token \citep{ouyang2022training}, that is,

\begin{equation}
    r_{t} = r_\varphi(q, o_{\le t}) - \beta \log\frac{\pi_{\theta}(o_{t}|q, o_{<t})}{\pi_{ref}(o_{t}|q, o_{<t})},
\label{eq:PPO-reward}
\end{equation}

where $r_\varphi$ designates the reward model, $\pi_{ref}$ refers to a reference policy—typically the starting SFT model—and $\beta$ controls the magnitude of the KL term.

Training a separate value function model in PPO, which is generally of similar scale to the policy, introduces substantial resource requirements in terms of both memory and compute. Additionally, within RL-based training, the value function serves as a baseline to reduce the variance in advantage estimation. In LLM settings, however, reward models tend to assign non-zero rewards only to the final token, making it challenging to train a value function that delivers accurate intermediate token predictions. To address these issues, and as depicted in Figure \ref{fig:grpo}, we introduce Group Relative Policy Optimization (GRPO). GRPO eliminates the necessity for a separate value function as in PPO; instead, it employs the average reward among multiple sampled outputs to the same prompt as a baseline. Concretely, given each question $q$, GRPO samples a group of outputs $\{o_1, o_2, \cdots, o_G\}$ using the previous policy $\pi_{\theta_{old}}$ and then updates the policy by optimizing the following loss:

\begin{equation}
\footnotesize
\begin{split}
    \mathcal{J}_{GRPO}(\theta) &= \mathbb{E}{[q \sim P(Q), \{o_i\}_{i=1}^G \sim \pi_{\theta_{old}}(O|q)]}  \\
    & \frac{1}{G}\sum_{i=1}^G\frac{1}{|o_i|} \sum_{t=1}^{|o_i|} \left\{ \min \left[ \frac{\pi_\theta(o_{i,t} | q, o_{i,<t})}{\pi_{\theta_{old}}(o_{i,t} | q, o_{i,<t})} \hat{A}_{i,t}, \text{clip} \left( \frac{\pi_\theta(o_{i,t} | q, o_{i,<t})}{\pi_{\theta_{old}}(o_{i,t} | q, o_{i,<t})}, 1 - \varepsilon, 1 + \varepsilon \right)  \hat{A}_{i,t} \right] - \beta \mathbb{D}_{KL}\left[\pi_{\theta} || \pi_{ref}\right]\right\} ,
\end{split}
\label{eq:GRPO-obj}
\end{equation}

where $\varepsilon$ and $\beta$ serve as hyper-parameters, and $\hat{A}_{i,t}$ indicates the advantage, computed solely from relative rewards among outputs sampled in the same group—a process that will be elaborated upon in subsequent subsections. GRPO’s method of estimating the advantages within each group closely matches the comparative structure of most reward model training data, which typically rely on pairwise comparisons of outputs in response to identical prompts. Moreover, rather than incorporating a KL penalty within the reward function as in PPO, GRPO directly includes the KL divergence between the current policy and the reference policy as a regularization term in the objective, which avoids additional complexity in computing $\hat{A}_{i,t}$. Distinct from the KL penalty used in (\ref{eq:PPO

\begin{algorithm}[t]
  \small
  \caption{Iterative Group Relative Policy Optimization}
  \textbf{Input} initial policy model $\pi_{\theta_{\text{init}}}$; reward models $r_\varphi$; task prompts $\mathcal{D}$; 
  hyperparameters $\varepsilon$, $\beta$, $\mu$
  \begin{algorithmic}[1]
    \State Set policy model $\pi_\theta \leftarrow \pi_{\theta_{\text{init}}}$
    \For{iteration = 1, \dots, I}
       \State Assign reference model $\pi_{ref} \leftarrow \pi_{\theta}$
      \For{step = 1, \dots, M}
      \State Randomly draw a batch $\mathcal{D}_b$ from $\mathcal{D}$
      \State Set the previous policy model $\pi_{\theta_{old}} \leftarrow \pi_{\theta}$
      \State For every question $q \in \mathcal{D}_b$, generate $G$ outputs $\{o_i\}_{i=1}^G$ by sampling from $\pi_{\theta_{old}} (\cdot \mid q) $
      \State For each sampled output $o_i$, obtain rewards $\{r_i\}_{i=1}^{G}$ using $r_{\varphi}$ 
      \State Determine $\hat{A}_{i,t}$ for the $t$-th token in $o_i$ via group relative advantage computation
      \For{GRPO iteration = 1, \dots, $\mu$}
        \State Adjust policy parameters $\pi_{\theta}$ to maximize the GRPO objective (see Equation \ref{eq:GC-GRPO})
      \EndFor
    \EndFor \State Improve $r_\varphi$ via continual training and replay
    \EndFor 
  \end{algorithmic}
  \textbf{Output} $\pi_\theta$
  \label{alg:iter-grpo}
\end{algorithm}

\subsubsection{Outcome Supervision RL with GRPO} More precisely, for each prompt $q$, a set of $G$ responses $\{o_1, o_2, \cdots, o_G\}$ is generated by the previous policy $\pi_{\theta_{old}}$. The reward model evaluates each response, assigning a corresponding set of scores $\mathbf{r} = \{r_1, r_2, \cdots, r_G\}$. These rewards are subsequently standardized by centering at the group mean and scaling by the group standard deviation. Outcome supervision supplies this normalized reward to each output $o_i$ after its completion and assigns the same normalized value to all token advantages $\hat{A}_{i, t}$ within $o_i$, formally, $\hat{A}_{i, t} = \widetilde{r}_i = \frac{r_i- {\rm mean}(\mathbf{r})}{{\rm std}(\mathbf{r})}$. The policy is then updated by maximizing the objective described in equation (\ref{eq:GRPO-obj}).

\subsubsection{Process Supervision RL with GRPO} Because outcome supervision limits feedback to the end of each output, it may offer insufficient learning signals for policy optimization in sophisticated mathematical contexts. Building upon prior work \cite{wang2023math}, we investigate process supervision, which supplies rewards at every intermediate reasoning step. Formally, given question $q$ and $G$ candidate outputs $\{o_1, o_2, \cdots, o_G\}$, the process reward model evaluates each intermediate step, producing a reward structure $\mathbf{R} = \{ \{r_1^{index(1)},\cdots,r_1^{index(K_1)}\}, \cdots,  \{r_G^{index(1)},\cdots,r_G^{index(K_G)}\} \}$, where $index(j)$ identifies the endpoint of the $j$-th step, and $K_i$ reflects the number of steps within output $o_i$. The rewards at each step are also normalized using their collective mean and standard deviation: $\widetilde{r}_i^{index(j)} = \frac{r_i^{index(j)} - {\rm mean(\mathbf{R})}}{{\rm std(\mathbf{R})}}$.
In this framework, the advantage for each token corresponds to the cumulative normalized rewards from subsequent steps, computed as $\hat{A}_{i, t} = \sum_{index(j) \ge t} \widetilde{r}_i^{index(j)}$. The policy optimization proceeds by maximizing the same objective function specified in equation (\ref{eq:

\subsubsection{Iterative RL with GRPO} As reinforcement learning training advances, the original reward model may become inadequate for effectively supervising the current policy model. Consequently, we investigate an iterative RL scheme utilizing GRPO. As depicted in Algorithm \ref{alg:iter-grpo}, this approach entails periodically regenerating training datasets for the reward model based on new samples drawn from the evolving policy model. The existing reward model undergoes continual training using a replay buffer containing 10\% previously collected data. Subsequently, the policy model is updated with reference to itself and further trained under the guidance of the refreshed reward model.

\subsection{Training and Evaluating DeepSeekMath-RL}

We implement RL on top of DeepSeekMath-Instruct 7B. For RL training, we use chain-of-thought-structured questions derived from GSM8K and MATH within the SFT data, yielding a dataset of approximately 144,000 problems. Other SFT datasets are intentionally excluded to examine RL’s effects on benchmark datasets without overlapping training data. Construction of the reward model’s training corpus follows the protocol of \citep{wang2023math}. The reward model is initially trained using DeepSeekMath-Base 7B, employing a learning rate of 2e-5. For GRPO, the policy model's learning rate is fixed at 1e-6, and we employ a KL penalty coefficient of 0.04. For every question, we generate $64$ sample completions. The maximum input length and batch size are set to 1024. The policy model undergoes a single parameter update at the conclusion of each exploration phase. DeepSeekMath-RL 7B is assessed using the same benchmarks as DeepSeekMath-Instruct 7B. Within this setup, GSM8K and MATH with chain-of-thought annotations constitute the in-domain evaluation, while all other benchmarks represent out-of-domain performance.

Table \ref{tab:sft_rl_math} presents results for both open- and closed-source models leveraging chain-of-thought and tool-augmented reasoning on English and Chinese tasks. Our analysis reveals the following insights: 1) DeepSeekMath-RL 7B achieves accuracy rates of 88.2\% on GSM8K and 51.7\% on MATH when employing chain-of-thought reasoning. These results exceed those of all publicly available models in the 7B–70B parameter range, and also outperform most proprietary models. 2) Notably, DeepSeekMath-RL 7B undergoes instruction tuning solely on chain-of-thought-formatted GSM8K and MATH data, building upon DeepSeekMath-Instruct 7B. Even with this limited training corpus, it demonstrates consistent improvements over DeepSeekMath-Instruct 7B on every metric assessed, underscoring the value added by reinforcement learning.

\section{Discussion}
Here, we detail key takeaways from our pre-training and RL experimental efforts.

\subsection{Lessons Learnt in Pre-Training}

We begin by describing our observations from pre-training. Unless explicitly mentioned otherwise, all training follows the procedures stated in Section \ref{sec:quality-policy}. In this context, the term DeepSeekMath Corpus specifically refers to an 89B-token dataset compiled during the second iteration of data acquisition.

\subsubsection{Code Training Benefits Mathematical Reasoning}

Although it has been commonly hypothesized—but not empirically substantiated—that training with code enhances reasoning abilities, we provide partial empirical evidence supporting this notion within the context of mathematical problem-solving: code training enhances mathematical reasoning capacities, both with and without access to external tools.

To investigate the impact of code pre-training on mathematical reasoning, we evaluated models using two approaches: a two-stage training strategy and a single-stage training setup.

\textbf{Two-Stage Training}

\begin{itemize}[topsep=0pt]
    \item \textbf{Code Training for 400B Tokens $\rightarrow$ Math Training for 150B Tokens}: DeepSeek-LLM 1.3B undergoes pretraining on 400B code tokens and subsequently on 150B math tokens;
    \item \textbf{General Training for 400B Tokens $\rightarrow$ Math Training for 150B Tokens}: For comparative purposes, we also conduct a training paradigm where, instead of code, 400B general tokens sourced from a comprehensive DeepSeek-AI dataset are employed in the initial phase, prior to math training, to examine the relative benefits of code versus general tokens for fostering mathematical reasoning.
\end{itemize}

\textbf{Single-Phase Training}

\begin{itemize}[topsep=0pt]
    \item \textbf{Math Training for 150B Tokens}: The model is exclusively trained for 150B tokens on math data;
    \item \textbf{Training on a mixture of 400B Code Tokens and 150B Math Tokens}: It is observed that subsequent math training after initial code training adversely affects performance on coding tasks. We therefore explore whether simultaneous exposure to both 400B code tokens and 150B math tokens in a unified stage can both bolster mathematical reasoning and curb catastrophic forgetting.
\end{itemize}

\paragraph{Results}
The downstream outcomes for each training regimen are presented in Table \ref{tab:code-math} and Table \ref{tab:code-math-general-eval}.

Code pretraining serves to enhance program-assisted mathematical reasoning, both in staged and combined training regimes. According to Table \ref{tab:code-math}, with a two-phase approach, the inclusion of code in the first phase already delivers a marked boost in tackling GSM8K and MATH benchmarks via Python execution. Subsequent math-focused training further elevates this performance. Moreover, in the unified training scenario, concurrently training on code and math tokens notably reduces the catastrophic forgetting characteristic of sequential phases, while simultaneously supporting both coding (Table \ref{tab:code-math-general-eval}) and program-facilitated mathematical reasoning (Table \ref{tab:code-math}).

Further, code training alone also supports improvements in mathematical reasoning without access to external tools. In the sequential framework, the preliminary code phase contributes moderate advancements and enhances the effectiveness of later math training, ultimately yielding superior outcomes. Conversely, when code and math data are intermixed during a single training phase, the ability to reason about mathematics without auxiliary tools is diminished. This may be attributed to the restricted capacity of DeepSeek-LLM 1.3B, which may not suffice for the concurrent assimilation of large volumes of code and math data.

\subsubsection{ArXiv Papers Appear Unhelpful for Advancing Mathematical Reasoning}

ArXiv papers have become a standard component in many mathematical pre-training datasets \citep{minerva,gpt-f,llemma,mathpile}, yet there is little comprehensive evidence regarding their influence on a model’s mathematical reasoning capabilities. Perhaps unexpectedly, our findings suggest that incorporating arXiv papers yields little benefit for this purpose. We assessed models at several scales, notably DeepSeek-LLM 1.3B and DeepSeek-Coder-Base-v1.5 7B \citep{deepseek-coder}, with two distinct arXiv-sourced corpora processed via different methodologies:

\begin{itemize}[topsep=0pt]
    \item \textbf{MathPile} \citep{mathpile}: a curated dataset comprising 8.9B tokens, assembled through specific cleaning and heuristic filtering, of which more than 85\% originates from scientific arXiv manuscripts;
    \item \textbf{ArXiv-RedPajama} \citep{redpajama}: a 28.0B-token dataset formed by stripping arXiv LaTeX documents of preambles, macros, comments, and references.
\end{itemize}

We trained DeepSeek-LLM 1.3B for 150B tokens and DeepSeek-Coder-Base-v1.5 7B for 40B tokens on each arXiv dataset independently. Our observations indicate that pre-training with only arXiv content neither enhanced the models’ performance on mathematical reasoning nor avoided regressions, across a variety of mathematical evaluation sets with varying difficulty. These assessments spanned tasks such as quantitative reasoning (GSM8K, MATH; see Table \ref{tab:arxiv-cot}), multiple-choice science and engineering problems (MMLU-STEM; Table \ref{tab:arxiv-cot}), as well as formal mathematics benchmarks (miniF2F; Table \ref{tab:arxiv_atp}).

Nonetheless, these findings are provisional, and should be interpreted cautiously. Notably, we have yet to investigate:

\begin{itemize}[topsep=0pt]
    \item The possible impact of arXiv data on specialized mathematical challenges outside this study, for example, the informalization of theorems, which involves translating formal statements or proofs to informal language;
    \item The consequences of mixing arXiv data with additional pre-training sources;
    \item Whether larger models might benefit more appreciably from arXiv-based pre-training.
\end{itemize}

These areas warrant further research, which we defer to subsequent investigations.

\subsection{Reflections on Reinforcement Learning}
\subsubsection{Toward a Unified Analytical Framework} This section offers a general analytical scheme to compare various training algorithms—such as SFT, RFT, DPO, PPO, and GRPO—while presenting experiments that examine key elements of this overarching structure. Typically, the gradient of a training objective with respect to model parameters $\theta$ is expressible in the form:

\begin{equation}
    \nabla_{\theta}\mathcal{J}_{\textcolor{red}{\mathcal{A}}}(\theta) = \mathbb{E}[\underbrace{(q,o) \sim \textcolor{red}{\mathcal{D}}}_{Data \ Source}]\left( \frac{1}{|o|} \sum_{t=1}^{|o|}  \underbrace{GC_{{\mathcal{A}}}(q, o, t, \textcolor{red}{\pi_{{rf}}})}_{Gradient \ Coefficient}  \nabla_{\theta}\log \pi_{\theta}(o_t | q, o_{<t})\right).
\label{eq:objective}
\end{equation}

This formulation consists of three principal elements: (1) the \textit{Data Source $\mathcal{D}$}, specifying the dataset used for training; (2) the \textit{Reward Function $\pi_{{rf}}$}, which provides the feedback signal for model updates; and (3) the \textit{Algorithm $\mathcal{A}$}, responsible for integrating both the training data and reward information into the gradient coefficient $GC$, thereby influencing the extent and direction of learning updates (either penalization or reinforcement). Within this unified framework, we consider and compare several prototypical training methodologies:

\begin{itemize}
    \item  \textbf{Supervised Fine-tuning (SFT)}: In SFT, a pretrained model undergoes additional supervised training using a curated dataset annotated by humans.
    \item \textbf{Rejection Sampling Fine-tuning (RFT)}: RFT refines an SFT-trained model further by selecting outputs generated in response to SFT prompts, retaining only those answers that pass correctness-based filtering for continued fine-tuning.
    \item \textbf{Direct Preference Optimization (DPO)}: DPO builds upon the SFT model, optimizing it by applying a pairwise preference-based objective to model outputs sampled from the SFT model, thereby emphasizing ranking according to human preferences.
    \item \textbf{Online Rejection Sampling Fine-tuning (Online RFT)}: In contrast to the standard RFT approach, Online RFT commences with an SFT-initialized model and updates it continuously, sampling outputs from the model itself in real time for augmentation and subsequent fine-tuning.
    \item \textbf{PPO/GRPO}: For PPO/GRPO, the initial policy is derived from the SFT model, and then iteratively improved by sampling outputs from the current (online) policy and reinforcing behaviors in line with the observed feedback.
\end{itemize}

A comprehensive summary of these approaches' elements can be found in Table \ref{tab:data_policy}. For an in-depth explanation of the derivation, please consult Appendix \ref{app:analysis-rl}.

\paragraph{Reflection on Data Source} We categorize the data sources utilized into two main types: online sampling and offline sampling. Online sampling refers to situations where the data used for training is gathered dynamically from the live policy model’s exploration, while offline sampling utilizes data drawn from the outputs of a pre-trained SFT model. Both RFT and DPO utilize offline-sampled data, whereas Online RFT and GRPO operate under the online sampling paradigm.

As illustrated in Figure \ref{fig:rl-analysis}, Online RFT exhibits notably superior performance compared to RFT across two evaluation benchmarks. More specifically, both methods perform comparably in the initial training phases, but as training progresses, Online RFT establishes a distinct performance lead, highlighting the benefits inherent in online data collection during training. This result is to be expected: during early training, the policy model closely resembles the SFT model, and thus the sampled data do not differ significantly. Over time, as the policy diverges, data sampled from the live policy reflect increasingly substantial differences, and the benefits of adaptive, real-time sampling become more pronounced.

\paragraph{Reflection on Gradient Coefficient} In the learning algorithm, the reward signal is incorporated into the computation of the gradient coefficient for parameter updates. In our experiments, we distinguish the reward function into two kinds: `Rule' and `Model.' The `Rule' setting assesses response quality by verifying the answer's correctness, whereas the `Model' approach involves training a reward model to generate scores for each response, where the training data for the reward model are themselves determined by rule-based judgments. Equations \ref{eq:GC-RFT} and \ref{eq:GC-GRPO} underscore an essential distinction between GRPO and Online RFT. GRPO uniquely modulates its gradient coefficient according to the magnitude of rewards given by the reward model, thus enabling both positive and negative reinforcement in proportion to reward intensity. In comparison, Online RFT does not implement such gradation: it fails to penalize incorrect outputs and applies uniform positive reinforcement to all correctly answered responses.

As illustrated in Figure \ref{fig:rl-analysis}, GRPO demonstrates better results than online RFT, emphasizing the effectiveness of modifying positive and negative gradient coefficients. Moreover, GRPO+PS outperforms GRPO+OS, underscoring the advantage of utilizing step-sensitive, fine-grained gradient coefficients. We further investigate iterative RL, performing two iterations in our experiments. Figure \ref{fig:iter-rl} reveals that iterative RL yields a notable boost in performance, with the most significant improvement occurring during the first iteration.

\subsubsection{Why RL Works?} This study applies reinforcement learning to a selected subset of instruction tuning data and observes a considerable performance increase over the baseline instruction tuning model. To better elucidate the effectiveness of reinforcement learning, we measure Pass@K and Maj@K accuracies for both the Instruct and RL-trained models across two benchmarks. Figure \ref{fig:maj-pass} shows that RL primarily lifts Maj@K performance without affecting Pass@K. These results suggest that RL strengthens overall model performance mainly by enhancing the robustness of the output distribution; that is, \textbf{the observed improvement primarily comes from promoting the correct answer within the TopK candidates, rather than directly enhancing core model competence.} Likewise, \citep{wang2023making} documented a \textbf{misalignment problem} in SFT model reasoning tasks, with additional research \citep{yuan2023rrhf,song2023preference,wang2023making} demonstrating that a series of preference alignment interventions can improve reasoning ability in such models.

\subsubsection{How to Achieve More Effective RL?}
\label{sec:effec_rl}
Our results show that reinforcement learning is highly effective on mathematical reasoning challenges. Furthermore, we introduce a unified conceptual framework for interpreting various representative training approaches, wherein all methods can be classified as either direct implementations or approximations of RL paradigms. As presented in Equation \ref{eq:objective}, three principal elements define these methods: Data Source, Algorithm, and Reward Function. We outline potential directions for future work with respect to each of these components.

\paragraph{Data Source} The data source forms the foundation for any training approach. In RL, we specifically regard the data source as comprising unlabeled questions together with outputs sampled from the policy model. Our experiments rely exclusively on questions from the instruction tuning phase, combined with basic nucleus sampling for output generation. We hypothesize that this methodological choice may explain why our RL method primarily benefits Maj@K accuracy. For future research, we intend to extend our RL framework to handle out-of-distribution prompts and to adopt \textbf{more sophisticated sampling (decoding) mechanisms}, such as those employing tree-search-based strategies \citep{yao2023tree}. In addition, leveraging \textbf{efficient inference techniques} \citep{xia-etal-2023-speculative,leviathan2023fast,kwon2023efficient,xia2024unlocking}, which determine how efficiently policy models can explore solution spaces, is expected to be of major significance.

\paragraph{Algorithms} Algorithms leverage both the data and the reward signal to compute the gradient, which is subsequently used to adjust the model parameters. In accordance with Equation \ref{eq:objective}, prevailing approaches operate under the assumption of complete \textbf{TRUST} in the reward function, employing it to either boost or suppress the conditional likelihood of specific tokens. Nevertheless, this reward signal cannot always be regarded as dependable, particularly when confronting highly intricate tasks. For instance, the PRM800K datasets \citep{lightman2023let}, despite being meticulously labeled by expert annotators, still exhibit an estimated 20\% rate of misannotations\footnote{\url{https://github.com/openai/prm800k/issues/12\#issuecomment-1728491852}}. Motivated by these observations, we aim to investigate reinforcement learning strategies that demonstrate resilience in the presence of noisy reward signals. We posit that \textbf{WEAK-TO-STRONG} \citep{burns2023weak} alignment methodologies hold the potential to fundamentally reshape learning algorithms.

\paragraph{Reward Function} The reward function furnishes the primary feedback for training. In reinforcement learning, it is often instantiated as a neural reward model. We identify three critical avenues for advancing reward models: 1) \textbf{Improving the generalization capacity of the reward model.} The reward model should generalize robustly to novel queries and high-quality decoding outputs beyond the training distribution; if not, reinforcement learning may only reinforce existing model behaviors instead of enhancing underlying competencies. 2) \textbf{Incorporating uncertainty estimation in reward models.} Accounting for model uncertainty could serve as a crucial connector between imprecise reward models and weak-to-strong learning frameworks. 3) \textbf{Efficient construction of precise process reward models} capable of delivering detailed feedback to support complex reasoning tasks \citep{lightman2023let,wang2023math}.

\section{Conclusion, Limitation, and Future Work} In this work, we introduce DeepSeekMath, which surpasses all open-source counterparts on the MATH benchmark at the competition level and narrows the gap with proprietary models. DeepSeekMath~builds on DeepSeek-Coder-v1.5 7B, having been further trained over 500B tokens, notably incorporating 120B math-specific tokens primarily collected from Common Crawl. Through comprehensive ablation analysis, we determine that web data can provide rich, high-quality mathematical resources, whereas data sourced from arXiv may not yield the anticipated benefits. We propose Group Relative Policy Optimization (GRPO), a modification of Proximal Policy Optimization (PPO), capable of enhancing mathematical reasoning performance while utilizing reduced memory resources. Our experiments indicate GRPO remains impactful, even as DeepSeekMath-Instruct 7B achieves strong benchmark results. Furthermore, we establish a unified framework for interpreting related methods and outline several avenues for advancing reinforcement learning approaches.

Despite the strong performance of DeepSeekMath~in quantitative reasoning tasks, it demonstrates limitations in geometry and theorem-proving relative to closed-source systems. For example, in our preliminary evaluations, the model struggles with triangle and ellipse-related problems, suggesting a potential bias in data selection during both pre-training and fine-tuning. Moreover, DeepSeekMath’s model size limits its performance in few-shot scenarios when compared to GPT-4: while GPT-4 benefits from few-shot examples, DeepSeekMath~exhibits similar outcomes under both zero-shot and few-shot settings. Going forward, we plan to refine our data selection methods to build a more robust and diverse pre-training dataset. We also aim to investigate further directions, as described in Section \ref{sec:effec_rl}, to promote more effective reinforcement learning strategies for large language models.